\documentclass[11pt,a4paper]{article}
\usepackage[margin=2.4cm]{geometry}
\usepackage{amsmath,amssymb}
\usepackage{booktabs}
\usepackage{xcolor}
\usepackage[colorlinks=true,linkcolor=blue,urlcolor=blue]{hyperref}

\newcommand{\ep}{\epsilon}
\newcommand{\gf}{\gamma_5}
\newcommand{\sla}[1]{#1\!\!\!/}
\newcommand{\tr}{\mathrm{Tr}}
\newcommand{\PS}{_{\rm PS}}
\newcommand{\note}[1]{\par\smallskip\noindent\colorbox{gray!12}{\parbox{\dimexpr\linewidth-2\fboxsep}{\small #1}}\par\smallskip}

\title{$\gamma_5$ in the Drell--Yan computation\\[2pt]
\large When is $aa=vv$? Larin's scheme, open lines and pure singlet}
\author{DY\_P2B --- notes\_Larin/}
\date{\today}

\begin{document}
\maketitle

\begin{abstract}
The vector boson $V=\gamma^*,Z$ couples to quarks through $\gamma^\mu(v_q+a_q\gf)$. Squaring gives vector--vector,
axial--axial and mixed pieces. We explain when the axial--axial piece equals the vector--vector one ($aa=vv$) and
when it does not, what Larin's prescription for $\gf$ does in squared amplitudes, and which axial-current
renormalisation constants are needed. The rule: $aa=vv$ when $V$ sits on the \emph{same open quark line} in the
amplitude and in the complex-conjugate amplitude, and only then. How to assemble all contributions with their
counterterms: \texttt{assembly.pdf}.
\end{abstract}

\section*{Summary}
\noindent Labels of the $V$ vertex (as in \texttt{assembly.pdf}): $v$, $a$ on open quark line 1 (flavour $q$); $v_2$,
$a_2$ on open quark line 2 (flavour $Q$); $v\PS$, $a\PS$ on a closed quark loop, the \textbf{pure singlet (PS)}, with
couplings $\sum_fv_f$, $\sum_fa_f$.
\begin{center}
\small
\begin{tabular}{lll}
\toprule
$V$ in amplitude / conjugate & example & axial--axial \\
\midrule
same open line & $v\,v^*$, $v_2v_2^*$ & $aa=vv$ (anticommuting $\gf$) \\
different open lines & $v\,v_2^*$ & $aa_2\neq vv_2$, also in $d=4$: Larin \\
identical quarks ($q=Q$) & both kinds, same couplings & split by line \\
closed loop (PS) & $v\PS v^*$, $a\PS a^*$ & vector: compute; axial: Larin, $\propto\sum_fa_f$ \\
\bottomrule
\end{tabular}
\end{center}
The parity-odd pieces ($va$, $av$) contain one $\gf$ in the hadronic trace and always need Larin
(section~\ref{sec:lept}).

\section{The leptonic tensor}\label{sec:lept}

For massless leptons, $V\to\ell^-(l_1)\ell^+(l_2)$ with vertex $\gamma^\mu(v_\ell+a_\ell\gf)$:
\begin{equation}
L^{\mu\nu}=\tr\big[\sla{l}_1\gamma^\mu(v_\ell+a_\ell\gf)\sla{l}_2\gamma^\nu(v_\ell+a_\ell\gf)\big]
=(v_\ell^2+a_\ell^2)\,4\big(l_1^\mu l_2^\nu+l_1^\nu l_2^\mu-g^{\mu\nu}l_1\!\cdot l_2\big)
+2v_\ell a_\ell\,\tr[\sla l_1\gamma^\mu\sla l_2\gamma^\nu\gf].
\end{equation}
(The conjugate vertex is again $\gamma^\nu(v_\ell+a_\ell\gf)$, since $\gamma^0(\gamma^\nu\gf)^\dagger\gamma^0=\gamma^\nu\gf$; the first
$(v_\ell+a_\ell\gf)$ moves through $\sla l_2\gamma^\nu$ without sign, and $(v_\ell+a_\ell\gf)^2=v_\ell^2+a_\ell^2+2v_\ell a_\ell\gf$.)
The symmetric part multiplies the parity-even hadronic pieces ($vv$, $aa$, \dots), the antisymmetric ($\epsilon$) part,
with coupling $2v_\ell a_\ell$, the parity-odd ones ($va$, $av$, \dots).

Averaged over the lepton angles at fixed $q=l_1+l_2$: $\langle l_1^\rho l_2^\sigma\rangle=A\,q^\rho q^\sigma+B\,q^2g^{\rho\sigma}$;
contracting with $g_{\rho\sigma}$ and $q_\rho q_\sigma$ gives $q^2/2=(A+4B)q^2$, $q^4/4=(A+B)q^4$, so $A=\tfrac16$, $B=\tfrac1{12}$:
the antisymmetric part averages to zero and
$\langle L^{\mu\nu}\rangle=\tfrac43(v_\ell^2+a_\ell^2)(q^\mu q^\nu-q^2g^{\mu\nu})=\tfrac43(v_\ell^2+a_\ell^2)L_1^{\mu\nu}$.
The parity-odd pieces contribute only to observables that resolve the lepton direction; there they are projected with
the antisymmetric projector $L_6$.

\section{\texorpdfstring{$\gamma_5$}{gamma5} in \texorpdfstring{$d$}{d} dimensions; Larin's scheme}\label{sec:g5}

In $d\neq4$ one cannot have both $\{\gf,\gamma^\mu\}=0$ and $\tr[\gf\gamma^\alpha\gamma^\beta\gamma^\gamma\gamma^\delta]\propto\epsilon^{\alpha\beta\gamma\delta}\neq0$
(anticommutation and cyclicity force that trace to vanish). Anticommuting $\gf$ is therefore unambiguous only if no
trace with an odd number of $\gf$ has to be computed.

Larin [hep-ph/9302240] replaces the axial vertex by
\begin{equation}
\gamma^\mu\gf\;\to\;\tfrac12(\gamma^\mu\gf-\gf\gamma^\mu)=\tfrac{i}{6}\,\epsilon^{\mu\nu\rho\sigma}\gamma_\nu\gamma_\rho\gamma_\sigma
\end{equation}
(normalisation of arXiv:2306.10170, eq.~(4.3)) and contracts the $\epsilon$ tensors in $d$ dimensions. No $\gf$ is ever
moved, so the scheme is consistent; the axial current must then be renormalised by
\begin{align}
Z_a^{ns}&=1-2C_Fa+\Big(\tfrac{11}{2}C_F^2-\tfrac{107}{36}C_FC_A+\tfrac1{18}C_Fn_f+\frac{C_F\beta_0}{\ep}\Big)a^2,\\
Z_s&=C_FT_R\Big(\frac{3}{2\ep}+\frac34\Big)a^2\quad\text{(singlet, per flavour)},
\end{align}
with $a=\alpha_s/(2\pi)$, $\beta_0=\tfrac{11}{6}C_A-\tfrac23T_Rn_f$. The finite parts restore the anticommuting-$\gf$ Ward
identities; the pole $C_F\beta_0/\ep$ is the $\overline{\rm MS}$ part. Conversion check (Larin's constants as I recall
them, in $a_s=\alpha_s/4\pi=a/2$, to be checked against the original): $Z_5^{ns}=1-4C_Fa_s+a_s^2(22C_F^2-\tfrac{107}9C_AC_F+\tfrac29C_Fn_f)$,
$Z^{ns}_{\overline{\rm MS}}=1+a_s^2(\tfrac{22}3C_AC_F-\tfrac43C_Fn_f)/\ep$; with $a_s^2=a^2/4$ these give the
expression above. The singlet finite constant has $\tfrac{31}{18}C_Fn_f$ instead of $\tfrac29C_Fn_f$: the difference
$\tfrac32C_Fn_fa_s^2=\tfrac34C_FT_Rn_fa^2$ is the $\tfrac34$ of $Z_s$ (arXiv:2306.10170, eqs.~(5.5), (5.6)).

\emph{Remark on arXiv:2306.10170, eq.~(5.13):} the $\bar G^{(0),n}$ coefficient is printed as $2\beta_0/\ep$; expanding
$Z_a^{ns}\sqrt{\bar a}$ gives $2C_F\beta_0/\ep$ ($C_F\beta_0/\ep$ from $(-2C_F)(-\beta_0/2\ep)$ and $C_F\beta_0/\ep$ from the
pole of $Z_a^{ns}$): probably a missing $C_F$.

\section{Same line: \texorpdfstring{$aa=vv$}{aa = vv}}\label{sec:same}

One open massless quark line, $V$ attached once in the amplitude and once in the conjugate, any number of gluons.
The spin-summed product is one trace,
\begin{equation}
\tr\big[\sla p\;\Gamma_1\,\gamma^\mu\gf\,\Gamma_2\;\sla p'\;\bar\Gamma'_2\,\gamma^\nu\gf\,\bar\Gamma'_1\big],
\end{equation}
$\Gamma_2$ = string from the $V$ vertex to the end of the line in the amplitude, $\bar\Gamma_2'$ the same in the
conjugate (reversed). Along a massless line propagators $\sla k$ and vertices $\gamma^\alpha$ alternate and the line ends
with a vertex, so $\Gamma_2=\sla k_1\gamma^{\alpha_1}\cdots\sla k_j\gamma^{\alpha_j}$ has $2j$ matrices; the same for
$\bar\Gamma_2'$. Between the two $\gf$ there are $2j+1+2j'+1$ matrices: an \emph{even} number. With anticommuting
$\gf$, moving one $\gf$ next to the other costs $+1$, and $\gf^2=1$:
\begin{equation}
\boxed{\;\tr[\cdots\gamma^\mu\gf\cdots\gamma^\nu\gf\cdots]=\tr[\cdots\gamma^\mu\cdots\gamma^\nu\cdots]\quad\Rightarrow\quad aa=vv\;}
\end{equation}
diagram by diagram, in any $d$, also with loops on the line (gluon loops, quark loops without $V$). No trace with a
single $\gf$ is computed: the two are removed before tracing. Helicity version: $P_\pm\gamma^\alpha=\gamma^\alpha P_\mp$,
$P_\pm\sla k=\sla kP_\mp$, so chirality $h$ is conserved along the line, $A_a(h)=h\,A_v(h)$ and
$\sum_h|A_a|^2=\sum_hh^2|A_v|^2=\sum_h|A_v|^2$. The parity-odd $\sum_hA_vA_a^*=\sum_hh|A_v|^2$ keeps one $\gf$ in the
trace: it needs Larin.

\paragraph{Larin breaks it at $O(\ep)$.} Larin's vertex is not $\gamma^\mu\gf$ in $d$ dimensions. Born:
\begin{equation}
vv=N_c(d-2)z^{-3},\qquad aa_{\rm Larin}=-\tfrac16N_c(d-2)(d-3)(d^2-8d+10)z^{-3},
\end{equation}
\begin{equation}
\frac{aa_{\rm Larin}}{vv}=(1-2\ep)\Big(1-\tfrac23\ep^2\Big).
\end{equation}
Equal at $d=4$, different at $O(\ep)$; the same pattern holds for the double-real same-line pieces. Larin
$+Z_a^{ns}$ and anticommuting $\gf$ are both consistent and give the same physical non-singlet result; the $O(\ep)$
differences cancel only in the full combination (reals, virtuals, collinear counterterms in the same scheme).

\note{\textbf{Choice.} Same-line parity-even pieces: anticommuting $\gf$, i.e.\ compute $vv$ and set $aa=vv$, at all
orders (RRR, RR, R, B and the non-singlet part of the loop contributions). Then $Z_a^{ns}$ is never needed for them.
Everything else with an axial vertex (cross-line, PS, parity-odd): Larin $+Z$'s.}

\section{Different lines: \texorpdfstring{$aa\neq vv$}{aa != vv}}\label{sec:cross}

$q\bar q\to VQ\bar Q$ with $V$ on the $q$ line in the amplitude and on the $Q$ line in the conjugate: for $q\neq Q$ the
product is two traces with one $\gf$ each, and the two $\gf$ cannot meet. With chiralities $h_1$, $h_2$ of the lines,
\begin{equation}
\sum_{h_1h_2}A_a^{(1)}A_a^{(2)*}=\sum_{h_1h_2}h_1h_2\,A_v^{(1)}A_v^{(2)*}\neq\sum_{h_1h_2}A_v^{(1)}A_v^{(2)*}:
\end{equation}
configurations with $h_1=-h_2$ enter with opposite sign, so $aa_2\neq vv_2$ already in four dimensions. Each trace has a
single $\gf$, so anticommuting $\gf$ is ambiguous: these pieces need Larin (and $Z_a^{ns}$ on both vertices in their
loop corrections). Their coupling is $a_qa_Q$; summed over $Q$, $\sum_Qa_Q$ cancels within complete generations
except for the $b$.

\paragraph{Identical quarks} ($q\bar q\to Vq\bar q$): both lines have the same couplings, so $v_q^2$ and $a_q^2$
multiply same-line products, cross-line products and the exchange interference together. $aa=vv$ holds only for the
same-line ones: the pieces must be separated by line, not by coupling.

\section{Closed quark loops: pure singlet}\label{sec:ps}

$V$ on a closed loop, summed over the $n_f$ flavours of the loop: couplings $\sum_fv_f$, $\sum_fa_f$; only in loop
amplitudes.
\begin{itemize}
\item \textbf{Vector PS.} Furry: zero with two gluons on the loop ($V$ + 1 parton at one loop; four-quark processes at
one loop; $q\bar q\to V$ at two loops). With three gluons on the loop ($d^{abc}$) it is non-zero: $V$ + 1 parton at
two loops, $q\bar q\to Vgg$ at one loop, $gg\to Vg$ at one loop. It does not cancel between flavours ($v_u\neq-v_d$).
\item \textbf{Axial PS.} One $\gf$ in the loop: the anomalous triangle; Larin $+Z_s$ is required. Example: the
one-loop $a\PS$ amplitude for $q\bar q\to Vg$ times the tree $a^*$ (an NNLO RV piece) is non-zero by itself
(Landau--Yang does not apply: one gluon is off shell); its coupling $a_q\sum_fa_f$ makes it vanish, together with its $Z_s$
counterterm, in the sum over a complete mass-degenerate doublet ($a_u=-a_d$).
\item With $n_f=5$ massless flavours the $b$ is left without its partner: $\sum_fa_f\neq0$. A correct treatment of
this remainder needs the massive top in the loop (arXiv:2306.10170, introduction); in DY it is usually small and
treated separately. To check in the N$^3$LO DY literature before deciding.
\item $gg\to Vg$: vector $\propto d^{abc}$, axial $\propto f^{abc}$, so their interference (parity-odd) vanishes by colour.
\item $\gamma^*$: $a_f=0$. $W$: no PS (a $W$ cannot couple to a loop that connects only to gluons).
\end{itemize}

\end{document}
